% SPDX-License-Identifier: GPL-3.0-or-later
% Copyright (C) 2026 Alessandro Greco
\section*{第 \ref{cap:01-anonimato} 章：匿名、监控与元数据}

\begin{soluzione}{es:01-metadati}
使用 HTTPS 时，互联网接入服务商能够看到网站的 IP 地址，有时还可以通过 DNS 解析以及 TLS 的 SNI 字段获得域名。它还可以观察时间特征、连接持续时间和流量大小。使用 VPN 时，接入服务商只能看到通向 VPN 服务器的连接，但 VPN 提供者能够看到此前接入服务商可能看到的全部信息。使用 Tor 时，接入服务商只能看到连接到了 Tor 中继或网桥，以及相应的时间特征和数据量。Tor 网络内部没有任何单个节点同时了解用户与目的地。
\end{soluzione}

\begin{soluzione}{es:01-termini}
使用虚构姓名的账户属于假名性。公告栏上一条没有署名的消息，在没有其他信息能识别作者的前提下属于匿名性。只写姓名首字母的信件属于较弱的假名性，往往仍可能识别作者。匿名现金捐款是匿名性的例子。
\end{soluzione}

\begin{soluzione}{es:01-insieme}
匿名集从五十名 Tor 用户缩小到其中会说意大利语、并且在凌晨三点仍然清醒的人，可能只剩寥寥数人。Tor 隐藏的是 IP 地址，而不是行为。语言、时间、写作风格和个人信息都可能在网络流量未泄露身份时，依然使用户被识别。
\end{soluzione}

\begin{soluzione}{es:01-vpn}
VPN 提供者同时知道用户身份和访问目的地。所有信任都集中在一个运营者身上；它可能记录活动、依法被要求交出数据，或者遭到入侵。Tor 把信任分散到三个独立运营的中继，任何一个都看不到通信两端。要关联用户与目的地，必须攻破或同时观察第一跳与最后一跳。
\end{soluzione}

\section*{第 \ref{cap:02-radici} 章：起源}

\begin{soluzione}{es:02-mix}
随机猜测时，正确概率为百分之一。如果观察者提供某批次一百条消息中的九十九条，就能识别自己送出的消息，并将唯一剩余的消息判定为受害者的消息。匿名性随之消失。这就是\emph{$n-1$ 攻击}。真实的混合网络通过更复杂的批处理、伪装流量以及对进入消息的控制来应对这种攻击。
\end{soluzione}

\begin{soluzione}{es:02-penet}
使用 Chaum 意义上的\emph{匿名回复地址}。用户为一串混合节点准备经过多层加密的\enquote{返回洋葱}，最内层包含用户地址及加密回复所需的密钥。回复者把消息附在洋葱上；每个混合节点剥去一层，但没有任何单个节点知道最终目的地。无须集中式路由表。
\end{soluzione}

\begin{soluzione}{es:02-crowds}
每一步以 $3/4$ 的概率把请求转发给另一位群体成员，以 $1/4$ 的概率发送到网站。经过的成员数量服从几何分布，平均为 $1/(1/4)=4$，并计入第一次转发。接收请求的成员无法知道上一位转发者是否也是原始发起者，因为中间转发也遵循同样规则。
\end{soluzione}

\section*{第 \ref{cap:03-nascita-tor} 章：从 NRL 洋葱路由到 Tor Project}

\begin{soluzione}{es:03-cinque}
在三跳结构中，恶意中间节点立刻就能知道通信两端所属的两个组织网络，因为入口和出口节点位于这些组织的网络内部。使用四跳时，攻破一端附近节点的人可以知道需要继续攻击哪一台节点才能重建路径；使用五跳则还需额外攻破一个节点。Tor 使用三跳，是因为中继不位于用户自己的网络，而且增加跳数会提高延迟，却无法大幅缓解入口与出口流量关联的主要威胁。守卫机制更适合解决这一问题。
\end{soluzione}

\begin{soluzione}{es:03-cipolla}
如果只使用一份通过节点长期公钥加密的洋葱，就没有前向保密。之后获得这些密钥的人可以解密以前保存的洋葱。逐跳扩展式构建则让每一跳执行临时密钥交换，使用后销毁临时秘密。客户端还可以一次扩展一跳，立即识别没有回应的中继并更换它。使用单个预先构造的洋葱，只有尝试完整路径后才会发现失败。
\end{soluzione}

\begin{soluzione}{es:03-fiducia}
这是开放式论证题。可以独立验证的因素包括公开的代码与规范、可复现构建、科学文献中的独立分析，以及由不同运营者运行的网络。仍然需要信任的方面包括：发布的二进制程序是否与公开源码一致（虽然可以部分验证）；目录权威节点是否串通；是否存在未知漏洞。应根据可验证性和分散信任评估安全性，而不是仅凭项目资金来源作出假设。
\end{soluzione}

\section*{第 \ref{cap:04-modello-minaccia} 章：威胁模型}

\begin{soluzione}{es:04-casi}
（a）可以：网络运营者只能看到连接到了 Tor 中继。（b）可以：网站看到的是出口中继的地址。（c）就身份匿名而言不可以：用户登录账户时主动标识自己。Tor 只能隐藏用户从哪里连接。（d）很可能不可以：同时观察用户端与报社端的攻击者，能够尝试进行时间关联，而这正是 Tor 威胁模型明确不保证抵御的情形。
\end{soluzione}

\begin{soluzione}{es:04-usabilita}
难用的系统只能吸引少数动机强烈的人，而这些人可能恰恰最需要保护。仅仅使用该系统，就足以使他们显得独特，留下一个很小的匿名集。易用的系统则可以吸引数百万普通用户，让有敏感需求的人更容易融入其中。
\end{soluzione}

\begin{soluzione}{es:04-modello}
这是开放式练习。潜在攻击者包括能够观察消息来源所在地网络的政府或雇主、电子邮件或消息服务提供商，以及控制通信两国之间网络的实体。Tor 可以阻止本地网络得知记者与谁通信，也防止服务端直接看到记者的 IP 地址。但如果没有端到端加密，Tor 无法保护消息内容；它也无法隐藏服务可见的账户和时间戳等元数据，不能抵御设备遭到攻破，亦不能防止同时观察两端的攻击者进行流量关联。使用 SecureDrop 等洋葱服务可以减轻其中若干风险。
\end{soluzione}

\section*{第 \ref{cap:05-crittografia} 章：密码学基础}

\begin{soluzione}{es:05-xor}
$C_1 \oplus C_2 = P_1 \oplus P_2$，因为重复使用的密钥流会相互抵消。若已知 $P_1$，就能恢复 $P_2=C_1\oplus C_2\oplus P_1$。
\end{soluzione}

\begin{soluzione}{es:05-shake}
16 字节输出正好等于 32 字节输出的前 16 字节：SHAKE 生成一条可以读取任意所需长度的输出流。因此，密钥派生函数可以申请较长的流，并将其切分为多个部分。然而，为不同用途请求不同长度\emph{并不能}实现域分离；仍然需要使用不同的域标签。
\end{soluzione}

\begin{soluzione}{es:05-dh}
两次计算出的共享秘密相同。全零公钥对应小阶点，不论私钥是什么，都可能产生全零的共享秘密，也就是所有人都知道的秘密。\texttt{cryptography} 库拒绝返回这样的结果。ntor 明确要求这种检查（\cref{sec:ntor}）；否则攻击者能够强制生成可预测的共享秘密。
\end{soluzione}

\begin{soluzione}{es:05-etichette}
没有域分离时，一种上下文中生成的值可能被另一个上下文错误接受。例如，用来认证握手消息的 MAC，可能被当作另一阶段的密钥或证明，也可能被使用相同函数与密钥的其他协议接受。域标签使不同用途的输入在设计上彼此区分。
\end{soluzione}

\section*{第 \ref{cap:06-architettura} 章：网络架构}

\begin{soluzione}{es:06-vista}
入口守卫：能看到用户 IP 地址、流量时间与数据量，但不知道目的地。中间中继：既不知道用户，也不知道目的地。出口中继：知道目的地，并可看到没有在应用层加密的内容，却不知道用户身份。HSDir：不会获得用户访问普通网站的信息。目录权威节点：通常不知道某人访问了哪个网站，但如果客户端\emph{直接}向权威节点获取目录文档，尤其在引导时，权威节点可能看到客户端 IP。如果客户端从后备目录或缓存获取文档，则看到请求的是这些服务器，而非权威节点。
\end{soluzione}

\begin{soluzione}{es:06-pr}
\texttt{Link=3-5} 表示支持最高到版本 5 的链路协议（\cref{cap:07-canali}）。\texttt{Relay=2-6} 包括 ntor、ntor-v3（\texttt{Relay=4}）和 CGO（\texttt{Relay=6}）。\texttt{FlowCtrl=1-2} 表示支持经过认证的 SENDME 与拥塞控制。如果中继同时具有 \texttt{Exit} 标志，并协商使用拥塞控制与 CGO，就可以充当 CGO 出口（\cref{sec:ntor3}）。
\end{soluzione}

\begin{soluzione}{es:06-livelli}
100\,KB 应用数据约需 $102400/498 \approx 206$ 个 DATA cell，链路上相应的数据量约为 $206 \times 514 \approx 106$\,KB。还要加上少量请求、流建立和 SENDME cell（大约七个，即每 31 个数据 cell 发送一个），以及 TLS 记录开销。不计初始握手，总体开销约比有效载荷高 5--7\%。
\end{soluzione}
